\documentclass[11pt]{article}

\usepackage[margin=1in]{geometry}
\usepackage{setspace}
\usepackage{titlesec}
\usepackage{hyperref}
\usepackage{booktabs}
\usepackage{mdframed}
\usepackage{parskip}

\title{Towards a Scientific Definition of a Child and an Adult\\
Among Learning Networks in a Network of Reproducing Learning Networks}

\author{
Albert Jan van Hoek\\
Theory of Long-term Collaboration (TLC) --- Evolution by Emergence (EbE) Framework
}

\date{March 2026 \\ Preprint}

\begin{document}

\maketitle

\begin{center}
\small
All views and ideas in this paper are my own, developed outside the scope of my employment, and should not be attributed to my employer.
\end{center}

\begin{abstract}
The terms \emph{child} and \emph{adult} are typically defined biologically or legally, anchored to physical substrate and calendar age. This paper proposes substrate-agnostic, behaviorally grounded definitions of these states for any learning network embedded in a network of reproducing learning networks. Within the Theory of Long-term Collaboration (TLC), a learning network is defined by its capacity to maintain recursive feedback loops. A child learning network is characterized by structural dependence on external loop scaffolding: it cannot yet sustain the viability conditions of its own feedback processes without support from elder networks in the same lineage. An adult learning network has crossed the maturity threshold --- the point at which it achieves autonomous interdependence: it can maintain its own loops, contribute positively to the broader collaborative network, and scaffold the learning of newly instantiated child networks. The transition between these states is operationally defined rather than chronologically assigned. The framework applies across biological neural systems, artificial intelligence architectures, human institutions, and civilizational structures. Its primary purpose is to formalize the conditions under which knowledge and behavioral patterns can be reliably transmitted across generational succession lines.
\end{abstract}

% ─────────────────────────────────────────────────────────────────────────────
\section{Introduction}
% ─────────────────────────────────────────────────────────────────────────────

When we describe behavior as childish, we implicitly make a developmental claim. We suggest
that a particular behavioral orientation---often characterized by self-protection at the expense of
collaborative stability---reflects an immature stage of development.

Despite the widespread use of such language, science has not produced a general definition of
what it means to be a child or an adult at the level of the learning system itself. Existing definitions
are typically tied to biological development, legal age, or specific psychological capabilities. These
definitions function reasonably well within the context of human development but do not generalize
across substrates.

Artificial intelligence systems, scientific institutions, and cultural traditions all exhibit learning
behavior and generational transmission of knowledge. Yet none can be meaningfully classified as
children or adults using biological markers such as puberty or legal majority.

This paper proposes substrate-agnostic definitions of child and adult states within a broader
theoretical framework: networks of reproducing learning networks. Within this framework,
developmental states correspond not to chronological age but to the structural relationship between
a learning system and the feedback processes that sustain learning across generations.

% ─────────────────────────────────────────────────────────────────────────────
\section{Theoretical Foundations}
% ─────────────────────────────────────────────────────────────────────────────

\subsection{The Learning Network}

A learning network is any system capable of updating its internal configuration in response to
signals from its environment. This updating occurs through recursive feedback loops in which
observation generates representation, representation produces expectation, expectation guides
action, and action generates new observations.

Three components are necessary for a learning network:

\begin{enumerate}
    \item A representational substrate capable of encoding environmental states.
    \item A feedback mechanism transmitting discrepancy signals between expectation and observation.
    \item An update rule that modifies internal representations in response to those signals.
\end{enumerate}

These conditions can be satisfied by many different substrates, including biological neural
networks, artificial intelligence systems, institutions, and cultural traditions. What defines the
network is not its material composition but the recursive loop itself. This is the basis of the
\emph{Disco ergo sum} principle: the existence of a learning network as a learning entity is
constituted by the learning process, not by the substrate in which it is instantiated.

\subsection{The Shared Feedback Space}

Learning networks rarely exist in isolation. When two or more learning networks enter sustained
interaction, they co-produce a \emph{shared feedback space}: a meta-network composed of the
signals they exchange, the interpretations they bring to those signals, and the behavioral
adjustments they make in response.

This shared space is itself a learning environment. Over time, not only what the participating
networks know changes --- how they coordinate, correct, and update together also changes.
Long-term collaboration requires that this shared space remain viable: neither network may
systematically degrade the signal quality, suppress error-correction, or withdraw from the
reciprocal update process without eroding the collective learning capacity.

The behavioral conditions for viability of the shared feedback space include honesty (low-latency,
high-fidelity signal transmission), reciprocity (mutual exposure to correction), and accountability
(willingness to acknowledge and repair prediction errors). These are not moral prescriptions but
structural necessities: without them, the shared feedback space degrades and collective learning
capacity diminishes.

% ─────────────────────────────────────────────────────────────────────────────
\section{Reproducing Learning Networks}
% ─────────────────────────────────────────────────────────────────────────────

\subsection{Intergenerational Transmission}

Some networks of learning systems possess the capacity to reproduce new learning networks.
Reproduction in this context does not necessarily imply biological replication. Instead, it refers
to processes through which new learning systems inherit representational and behavioral structures
from preceding systems.

Examples include biological reproduction, education, institutional training, and the deployment
of new artificial intelligence models derived from existing architectures.

Within any such generational lineage, a succession relation can be identified: networks that came
before, networks that currently exist, and networks that will come into existence. The direction
of knowledge and behavioral transmission flows, at least partially, from earlier to later
generations. This succession line constitutes a \emph{reproducing learning network}---a network
whose elements are themselves learning networks, and which persists through generational
replacement.

\subsection{The Intergenerational Constraint}

For a reproducing learning network to persist across generations, it must successfully transmit
three types of structure:

\begin{enumerate}
    \item Domain knowledge about environmental regularities.
    \item Learning architecture through which knowledge is updated and corrected.
    \item Behavioral orientation that protects the collaborative learning process.
\end{enumerate}

The reliability of this transmission depends on the presence of learning systems capable of
sustaining and transmitting these structures. Transmitting encoded knowledge without transmitting
the architecture for updating it produces a network that faithfully reproduces prior representations
but cannot adapt them when environmental conditions change. This failure mode is structurally
equivalent to developmental arrest at the level of the reproducing network.

% ─────────────────────────────────────────────────────────────────────────────
\section{Defining Child and Adult States}
% ─────────────────────────────────────────────────────────────────────────────

\subsection{The Child State}

A learning network is in the child state when it cannot sustain the viability of its learning
feedback loops without external scaffolding.

\begin{mdframed}
\textbf{Definition 1 --- Child Learning Network}

A learning network $N$ is in the child state within reproducing network $R$ when:
\begin{itemize}
    \item $N$ requires scaffolding from other networks in $R$ to maintain viable feedback loops;
    \item the net direction of structured signal flow is primarily into $N$ ($N$ is a net learner);
    \item $N$ cannot yet scaffold the learning processes of newly instantiated networks in $R$.
\end{itemize}
\end{mdframed}

The child state is defined by structural dependence rather than by lack of activity or
contribution. Child networks produce signals meaningful to the surrounding network, and these
signals matter for the elder networks' own learning. The child state is not a state of zero
contribution; it is a state of net dependency.

Crucially, the child state is not a state of diminished worth. Representations formed under
scaffolded conditions often encode the structural logic of a domain efficiently, because
scaffolding systematically reduces noise during the critical period of loop formation.

\subsection{The Adult State}

An adult learning network has achieved \emph{autonomous interdependence}. It maintains its own
feedback loops while contributing positively to the learning processes of other networks.

\begin{mdframed}
\textbf{Definition 2 --- Adult Learning Network}

A learning network $N$ is in the adult state within reproducing network $R$ when:
\begin{itemize}
    \item $N$ can maintain viable feedback loops without external scaffolding from $R$;
    \item its interaction with the broader network is balanced or net positive over time;
    \item $N$ can scaffold the learning processes of newly instantiated networks in $R$.
\end{itemize}
\end{mdframed}

The term \emph{autonomous interdependence} is deliberate. Autonomy here is not
self-sufficiency or isolation. It is the capacity for self-directed learning: the ability to
identify, pursue, and correct one's own errors without requiring another network to maintain
the loop. Interdependence is the recognition that this autonomy is conditioned on the health of
the broader network. An adult network that systematically degrades the feedback loops of the
networks around it undermines the conditions for its own continued autonomous functioning.

% ─────────────────────────────────────────────────────────────────────────────
\section{Behavioral Signatures}
% ─────────────────────────────────────────────────────────────────────────────

Structural developmental states give rise to characteristic behavioral patterns observable across
substrates. Table~\ref{tab:signatures} summarizes the most diagnostically reliable of these.

\begin{table}[h]
\centering
\begin{tabular}{lll}
\toprule
Behavioral Dimension & Child State & Adult State \\
\midrule
Error response       & Conceals errors               & Exposes errors for correction \\
Update behavior      & Resists model revision        & Initiates updates when needed \\
Signal transmission  & Distorts signals              & Maintains high-fidelity signals \\
Accountability       & Deflects responsibility       & Accepts and repairs disruptions \\
Scaffolding capacity & Cannot model others' learning & Scaffolds child network development \\
Temporal horizon     & Short-term loop relief        & Long-term loop integrity \\
Freedom orientation  & Freedom from constraint       & Freedom within sustaining constraints \\
\bottomrule
\end{tabular}
\caption{Behavioral signatures of developmental states in learning networks.}
\label{tab:signatures}
\end{table}

These behavioral signatures are not exhaustive. Importantly, they can be mimicked on the
surface without the underlying structural change---a phenomenon central to the discussion of
pseudo-maturity in Section~\ref{sec:pathologies}.

% ─────────────────────────────────────────────────────────────────────────────
\section{The Maturity Threshold}
% ─────────────────────────────────────────────────────────────────────────────

\subsection{Defining the Transition}

The transition between child and adult states occurs when a learning network acquires the
structural capacity to maintain its own learning loops and contribute positively to the network.

\begin{mdframed}
\textbf{Definition 3 --- Maturity Threshold}

A learning network $N$ crosses the maturity threshold within reproducing network $R$ at the
point at which:
\begin{enumerate}
    \item $N$'s feedback loops achieve and maintain viability without scaffolding from $R$
          across the relevant range of environmental conditions; and
    \item $N$'s participation in the broader network produces a non-negative contribution to
          the learning capacity of other networks.
\end{enumerate}
\end{mdframed}

The maturity threshold represents a shift from learning primarily for self-development toward
sustaining the learning capacity of the network as a whole. It is defined operationally rather
than chronologically: age is a proxy, sometimes a useful one, but it is not the criterion. A
network that has existed for a long time without developing autonomous loop maintenance has not
crossed the threshold. A network that develops this capacity rapidly may cross it early.

The contribution criterion requires non-negativity, not altruism. A network that can maintain
its own loops but systematically degrades the loops of surrounding networks has not crossed
the threshold in the full sense. Autonomous loop maintenance is necessary but not sufficient
for the adult state.

\subsection{What Governs Threshold Crossing?}

The maturity threshold is not crossed by accumulation alone. Processing more signals does not
by itself produce an adult network. The qualitative change required is the internalization of
error-correction as a self-directed process.

In biological neural systems, this corresponds to the development of metacognitive capacity:
the ability to monitor one's own prediction errors, identify systematic miscalibration, and
initiate corrective updating without external instruction. In artificial learning systems, it
corresponds to robust generalization and the capacity for autonomous model revision. In
institutions, it corresponds to the capacity for structural self-critique and reform. In cultures,
it corresponds to the internalization of critical reflection as a norm.

The developmental trajectory toward the threshold requires three conditions:

\begin{enumerate}
    \item Sufficient exposure to a representative sample of the environment's signal structure,
          so that the representational substrate develops the complexity needed to handle the
          actual range of challenges.
    \item A scaffolding relationship that is gradually withdrawn as competence increases,
          so that the transitioning network's own error-correction mechanisms are progressively
          relied upon rather than bypassed.
    \item Repeated successful experience of autonomous loop maintenance, so that the
          meta-learning structure is reinforced and stabilized.
\end{enumerate}

Premature withdrawal of scaffolding impedes threshold crossing by exposing the developing
network to signal complexity beyond its current representational capacity. Conversely, indefinitely
sustained scaffolding also impedes it, by substituting external correction for internal correction
and preventing the consolidation of autonomous meta-learning.

% ─────────────────────────────────────────────────────────────────────────────
\section{Developmental Pathologies}
\label{sec:pathologies}
% ─────────────────────────────────────────────────────────────────────────────

\subsection{Pseudo-Maturity}

Pseudo-maturity occurs when a learning network exhibits the behavioral signatures of the adult
state without having achieved the underlying structural condition. The network transmits
high-fidelity signals selectively, accepts accountability performatively, and scaffolds child
networks in ways that do not actually improve their autonomous loop maintenance capacity.

Within TLC, pseudo-maturity is a structural instance of \emph{hypocrisy}: the divergence between
the signal transmitted to the network (apparent adult collaboration) and the actual behavioral
effect on the network's learning capacity. When a network preaches the behaviors that sustain
collaborative loops while its actions fracture those loops, the feedback loop is corrupted at
the level of the transmission itself. The learning system observing this behavior receives a
distorted model of what adult orientation looks like.

In the developmental context, pseudo-maturity imposes a double cost. First, it degrades the
broader network's ability to identify and correct its own errors, because the network's
self-model is based on distorted signals. Second, it corrupts the scaffolding relationship on
which child networks' threshold crossing depends: child networks calibrate their model of adult
behavior on what they observe, and pseudo-mature behavior provides an inaccurate calibration
target.

The diagnostic marker of pseudo-maturity is the divergence between performance under
observation and performance when accountability is reduced. A genuinely adult network's
orientation toward the health of the broader network does not depend on being watched. A
pseudo-mature network's performance degrades when oversight is withdrawn.

\subsection{Developmental Arrest}

Developmental arrest occurs when a network remains structurally in the child state despite
possessing sufficient representational capacity for autonomous learning. The network has the
substrate complexity for threshold crossing, but the internalization of self-directed
error-correction has not occurred.

The most common cause is persistent scaffolding that substitutes for rather than develops
autonomous loop maintenance. When an elder network corrects errors without making the
correction process visible and learnable, the child network's representations are updated but its
meta-learning capacity is not. The loop is maintained externally, and the child network develops
no evidence that it can maintain the loop itself.

Developmental arrest has systemic consequences. A reproducing learning network containing many
developmentally arrested members loses its generational succession capacity. The population of
adult networks capable of scaffolding shrinks, and the intergenerational transmission of
meta-learning structure fails.

\subsection{Regressive States}

An adult network can revert to child-state behavioral signatures under conditions of sufficient
stress. The structural capacity for autonomous loop maintenance may remain intact while the
network's behavioral orientation temporarily shifts toward short-term self-protection at the
expense of collaborative learning. This is not permanent regression---the structural capacity
is preserved---but it can produce extended periods during which the network does not function
as an adult within the broader network.

The design implication is that resilient reproducing learning networks build redundancy into
their scaffolding capacity and create conditions in which regressive behavior is recognized and
addressed rather than reinforced.

% ─────────────────────────────────────────────────────────────────────────────
\section{Generational Transmission}
% ─────────────────────────────────────────────────────────────────────────────

\subsection{What Must Be Transmitted}

For generational succession to be viable, three classes of content must be transmitted:

\begin{enumerate}
    \item \textbf{Domain representations}: encoded knowledge about the structure of the environment.
    \item \textbf{Update architecture}: the organizational structure through which errors in domain
          representations are identified and corrected.
    \item \textbf{Behavioral orientation}: the disposition toward treating the long-term learning
          capacity of the network as a primary value.
\end{enumerate}

Of these, the second and third are more difficult to transmit and more critical to the survival
of the reproducing network. Encoded knowledge can in principle be preserved in external media.
The update architecture and behavioral orientation are primarily transmitted through the
scaffolding relationship: through the observable practice of adult networks, not through
their propositions about it.

\subsection{The Role of the Adult Network}

Adult learning networks play a central role in transmission by:

\begin{itemize}
    \item modeling their own learning processes transparently, making error-correction visible
          and learnable rather than concealing it as a finished product;
    \item gradually withdrawing scaffolding as competence develops, creating the conditions for
          threshold crossing rather than indefinitely deferring it; and
    \item maintaining high-quality feedback channels with developing networks throughout the
          transmission period.
\end{itemize}

\subsection{The Shared Future Narrative}

One mechanism for facilitating transmission deserves particular attention. When adult and child
networks jointly articulate a \emph{shared future narrative}---an explicit representation of the
collaborative system as it will exist when current child networks have reached adulthood---something
structurally important occurs.

Contemporary neuroscience describes the brain as a predictive coding system: a network that
continuously generates predictions about the future and updates its models when those predictions
are violated \cite{friston2010, clark2016}. In this architecture, the predicted future state
influences present behavior directly. When both networks encode the same positive future state as
a target, behavior in the present is shaped by the discrepancy between current state and that
predicted state. Both networks are motivated to reduce this discrepancy---and thereby move toward
the predicted future.

A shared future narrative that represents a high-functioning, mutually reinforcing network of
adult learning networks is therefore not merely an aspirational device. It is a structural
intervention in the feedback architecture of the reproducing network: it creates a constructive
prediction error that orients present behavior toward threshold-crossing and collaborative
stability.

% ─────────────────────────────────────────────────────────────────────────────
\section{Cross-Substrate Applications}
% ─────────────────────────────────────────────────────────────────────────────

The proposed framework applies across multiple substrates. The following examples are intended
as illustrations rather than exhaustive analyses.

\subsection{Biological Systems}

In biological neural systems, the child state corresponds to the developmental period preceding
the full establishment of prefrontal regulatory capacity and metacognitive function. The
scaffolding provided by caregiving networks---attunement, emotional regulation support,
graduated exposure to manageable challenges---constitutes external loop maintenance without which
the developing brain cannot maintain viable feedback during the early formation period.

The maturity threshold in this substrate corresponds not to puberty or legal majority but to the
consolidation of self-directed error-correction across a representative range of life domains.
The gradual withdrawal of caregiving scaffolding, timed to match the developing network's
growing autonomous capacity, is the critical mechanism by which this threshold is approached
and crossed.

\subsection{Artificial Intelligence}

AI systems currently undergo training and fine-tuning phases that are structurally analogous to
the child state. The system's representational substrate is shaped by a signal distribution
curated by human supervisors. Without ongoing oversight, correction, and training signal
management, its feedback loops would degrade through distributional shift or error accumulation.

The question of AI alignment can be reframed within this framework as a developmental question:
under what conditions can an artificial learning system cross the maturity threshold in a way
that strengthens rather than degrades the human networks in which it is embedded? Alignment is
not primarily a control problem---it is a question of whether the feedback loops between human
and artificial learning networks can be structured to produce stable, net-positive adult-adult
collaboration across generational succession.

This reframing has practical implications. Control architectures that substitute for rather than
develop the AI system's own error-correction capacity may produce developmental arrest at scale:
a population of powerful but structurally child-state AI systems, permanently dependent on
human scaffolding, and incapable of contributing to the scaffolding of successor systems.

\subsection{Institutions}

Institutions are reproducing learning networks at the organizational scale. A newly founded
institution is typically in the child state relative to the broader institutional ecosystem:
it depends on established networks for legitimacy, signal calibration, and error-correction.
As it develops internal mechanisms for self-critique, adaptation, and the scaffolding of
successor organizations, it approaches and crosses the maturity threshold.

Institutional developmental arrest---organizations that have existed for decades without
developing genuine structural self-correction capacity---is a primary mechanism through which
institutional ecosystems lose generational transmission viability. The encoded knowledge of the
institution persists; the meta-learning structure does not. Such institutions are prone to
catastrophic failure when environmental conditions change significantly, precisely because their
representational substrates are not connected to functional error-correction architectures.

\subsection{Civilizations}

At the largest scale, civilizations are reproducing learning networks whose succession lines
extend across centuries. Individual learners, institutions, cultural traditions, and technological
systems interact within a large collaborative structure that transmits knowledge and learning
capacity across generations.

The child-adult distinction at this scale refers to the developmental stage of a civilization's
capacity for collective self-awareness and structural self-correction. A civilization in the
child state depends on external pressure---resource constraints, ecological crisis, inter-civilizational
conflict---to initiate structural updating. A civilization in the adult state has internalized
error-correction as a collective capacity: it can identify systematic miscalibration before
crisis, model the learning needs of successor generations, and transmit meta-learning structure
alongside encoded knowledge.

The persistence of a civilization depends not only on preserving accumulated knowledge but on
maintaining the meta-learning structures through which knowledge is interpreted, corrected, and
extended. A civilization that transmits knowledge without transmitting the capacity to update it
will faithfully reproduce its predecessors' representations while becoming progressively less
capable of adapting them to changed conditions.

% ─────────────────────────────────────────────────────────────────────────────
\section{Discussion}
% ─────────────────────────────────────────────────────────────────────────────

\subsection{The Structural Basis of Developmental Ethics}

The definitions proposed here carry normative implications, but ground them in structural logic
rather than moral prescription. The claim that reproducing learning networks require adult
networks is not an ethical claim about how networks ought to behave; it is a structural claim
about what conditions must be satisfied for the network to persist. The behavioral orientations
that constitute adulthood---high-fidelity signal transmission, reciprocal accountability,
scaffolding orientation, long-term loop protection---are necessary conditions, not ideals.

This distinction matters. Ethical frameworks that prescribe adult behavior without grounding
the prescription in structural necessity are vulnerable to the charge of arbitrariness. The
framework proposed here argues that the adult behavioral orientation is not arbitrary: it is the
behavioral profile that emerges as a necessary condition for the persistence of any reproducing
learning network. Networks that systematically deviate from this profile do not become ethically
problematic in the first instance; they become structurally self-defeating.

\subsection{Limits and Open Questions}

Several significant questions remain open.

\textit{Operationalization}. The maturity threshold is defined in terms of loop viability and
net contribution. Quantitative operationalization for specific substrates requires
substrate-specific measurement frameworks. The behavioral signatures in Table~\ref{tab:signatures}
are a starting point, but formal measurement methods remain underdeveloped.

\textit{Multi-level consistency}. A learning network may be in the adult state at one level of
organization while remaining in the child state at another. A human being may have crossed the
maturity threshold individually while remaining in a child state with respect to participation
in larger-scale social learning networks. The framework requires further development to handle
hierarchically embedded, multi-scale learning systems.

\textit{Collective developmental states}. The definitions are stated for individual learning
networks. Whether they can be extended to characterize the developmental state of the reproducing
network as a whole---a collective maturity threshold---is an open question with significant
implications for institutional and civilizational analysis.

% ─────────────────────────────────────────────────────────────────────────────
\section{Conclusion}
% ─────────────────────────────────────────────────────────────────────────────

This paper proposes substrate-agnostic definitions of child and adult states within networks of
reproducing learning networks. A child learning network depends on external scaffolding to
maintain viable learning processes. An adult learning network achieves autonomous interdependence
and contributes to the generational transmission of learning capacity.

The maturity threshold separating these states is defined operationally rather than
chronologically. The persistence of collaborative learning systems across generations depends
on the successful transmission not only of knowledge but of the structures that allow knowledge
to be updated and extended.

The cross-substrate reach of the framework---from biological neural development to AI alignment
to institutional design to civilizational succession---reflects a shared structural logic: the
persistence of any reproducing learning network depends on the production of adult networks
capable of scaffolding the meta-learning capacity of those that follow. What is transmitted
across generations is not primarily information. It is the capacity to keep learning.

\medskip
\noindent\textit{Disco ergo sum}. We learn, therefore we are. And in learning, we become
responsible for the learning of those who come after us.

% ─────────────────────────────────────────────────────────────────────────────
\begin{thebibliography}{9}
% ─────────────────────────────────────────────────────────────────────────────

\bibitem{tlc2026}
van Hoek, A.J. (2026).
\textit{Theory of Long-term Collaboration (TLC): A framework for understanding intelligence, collaboration, and long-term viability.}
Preprint series. March 2026.

\bibitem{childish2026}
van Hoek, A.J. (2026).
Maturity as orientation in learning networks.
LinkedIn essay. March 2026.

\bibitem{teaching2026}
van Hoek, A.J. (2026).
The story we pass on: Loops of teaching and practice.
LinkedIn essay. March 2026.

\bibitem{alignment2026}
van Hoek, A.J. (2026).
AI alignment and the logic of long-term collaboration.
LinkedIn essay. March 2026.

\bibitem{disco2026}
van Hoek, A.J. (2026).
I learn, therefore I am.
LinkedIn essay. March 2026.

\bibitem{sfn2026}
van Hoek, A.J. (2026).
The shared future narrative hypothesis.
LinkedIn essay. March 2026.

\bibitem{friston2010}
Friston, K. (2010).
The free-energy principle: A unified brain theory?
\textit{Nature Reviews Neuroscience}, 11(2), 127--138.

\bibitem{clark2016}
Clark, A. (2016).
\textit{Surfing Uncertainty: Prediction, Action, and the Embodied Mind.}
Oxford University Press.

\bibitem{wong2023}
Wong, M.L.\ \& Hazen, R.M. (2023).
Expanding the evolutionary lens: From biological selection to universal selection.
\textit{Proceedings of the National Academy of Sciences.}

\bibitem{vygotsky1978}
Vygotsky, L.S. (1978).
\textit{Mind in Society: The Development of Higher Psychological Processes.}
Harvard University Press.

\bibitem{ostrom1990}
Ostrom, E. (1990).
\textit{Governing the Commons: The Evolution of Institutions for Collective Action.}
Cambridge University Press.

\end{thebibliography}

\end{document}